
Era un conjunto de parches del kernel de Linux que trataban de convertirlo en un \textit{RTOS} mediante dos enfoques distintos. Uno era el uso de un microkernel, siguiendo la estrategia de \textit{RTAI}. El otro era modificar el código del kernel en sí mismo. Con el paso del tiempo, gran parte de los parches fueron reescritos, y en 2015, comenzaron esfuerzos desde la \textit{Linux Foundation} de comenzar el \textit{Real-Time Linux Collaborative Project}. Dicha iniciativa, posiciona el proyecto donde está ahora, recientemente fusionado con la \textit{mainline} de Linux el 30 de Septiembre de 2024 \cite{RTLdeveloperteam2022, RTLdeveloperteam2024}.

Es por esto por lo que, \textit{Preempt-RT}, es la opción más atractiva para apostar en un futuro, pues al contar con el apoyo de la \textit{Linux Foundation} cuenta con soporte oficial y esto hace que tenga mayor potencial de mejora.
